% ===================================================================
%  TABEL Nr. 16 — Die groot warenhuis. — Die basaar. — Die speelgoed.
%  Meme regime que les precedents. Tableau a UNE SEULE COLONNE, avec
%  une note en pied placee entre le 8 et le 9.
%
%  LE BLOC t16-15-1 EST L'UN DES TROIS ECARTS DECLARES, et pour une
%  raison de COMPOSITION : le fac-simile ido y ecrit « d) » sans sa
%  parenthese ouvrante, si bien que le releve des renvois ne le voit
%  pas. Les traductions ecrivent le renvoi entier — le neerlandais et
%  le tcheque le font — et cette colonne aussi.
%
%  CE MEME BLOC EST LE SEUL DU LIVRET A NUMEROTER PAR LETTRES : de
%  (a) a (f), pour la table d'astronomie. Les lettres se cliquent sur
%  la page de lecture comme les numeros.
%
%  ET LE 5 EST UN COURS DE STRATEGIE EN TROIS ALINEAS ou presque
%  chaque mot est en gras sans porter de renvoi. Ce ne sont pas des
%  termes de planche, ce sont les mots que le grand-pere apprend a
%  l'enfant. On garde le gras : il est dans le fac-simile.
%
%  DERNIER TABLEAU, ET DERNIER PIEGE NEERLANDAIS : neushoorn,
%  struisvogel, luipaard, schildpad s'ecrivent presque seuls. Ce sont
%  renoster, volstruis, luiperd, skilpad.
% ===================================================================

%%K t16-apar-1 apar
\VUsaut{4.0mm}
\VUtitre{15.0pt}{60}{TABEL Nr. 16}
\VUsaut{2.0mm}
\VUpk{13.2pt}{40}{Die groot warenhuis. --- Die basaar.}
\VUpk{13.2pt}{40}{Die speelgoed.}
\VUsaut{3.4mm}

%%K t16-01-1 p
1. --- Die nuwe jaar kom nader. Die groot warenhuise het hulle uitstallings van
\VUgras{speelgoed} \textsuperscript{(1)} en nuwejaarsgeskenke gereedgemaak.

%%K t16-01-2 p
Ek het my oupa gevra of hy my na die basaar wou bring om daardie mooi uitstalling te
sien, en hy het daarin toegestem.

%%K t16-02-1 p
2. --- Eers het ons voor mooi wapenskilde stilgehou wat 'n \VUgras{geweer}, 'n
\VUgras{kepie}, 'n \VUgras{sabel}, 'n \VUgras{degengordel} en 'n paar
\VUgras{epoulette} bevat, of 'n \VUgras{helm}, 'n \VUgras{kuras} \textsuperscript{(3)}
en 'n \VUgras{pistool} \textsuperscript{(4)}. Dit alles het my baie meer geboei as die
\VUgras{towerlanterns} \textsuperscript{(2)}.

%%K t16-tit-1 sub
\VUsaut{3.0mm}
\VUpk{10.2pt}{40}{Mooi gesigte.}
\VUsaut{2.6mm}

%%K t16-03-1 p
3. --- In dieselfde afdeling was 'n \VUgras{skildwag} \textsuperscript{(5)} in sy
\VUgras{skilderhuisie}; hy het gelyk of hy wag hou voor 'n hele menasserie van karton.
Hoeveel diere was daar: 'n \VUgras{Papegaai} \textsuperscript{(6)} op sy stok, 'n
\VUgras{renoster} \textsuperscript{(7)} met 'n horing op die neus, 'n \VUgras{rendier}
\textsuperscript{(8)}, 'n \VUgras{volstruis} \textsuperscript{(9)}, 'n \VUgras{aap}
\textsuperscript{(10)} wat 'n neut kraak, 'n \VUgras{jakkals} \textsuperscript{(11)} wat
'n hen wegdra, 'n \VUgras{krokodil} \textsuperscript{(12)} wat met enorme kake gaap, 'n
\VUgras{kameel} \textsuperscript{(13)}, 'n elegante \VUgras{windhond}
\textsuperscript{(14)}, 'n \VUgras{panter} \textsuperscript{(15)}, daarna twee diere wat
ek nie geken het nie en wat, na wat gesê word, 'n \VUgras{wesel} \textsuperscript{(16)}
en 'n \VUgras{hiëna} \textsuperscript{(17)} is. Ook was daar 'n \VUgras{luiperd}
\textsuperscript{(18)} en 'n \VUgras{wolf} \textsuperscript{(21)}; heel naby was
lieflike \VUgras{rotjies} \textsuperscript{(19)} en piepklein \VUgras{muisies}
\textsuperscript{(20)} van rubber. Laastens het 'n \VUgras{seekoei}
\textsuperscript{(22)} en 'n boggelrige \VUgras{dromedaris} \textsuperscript{(23)} die
versameling afgesluit. Ek sou daar nog meer ure gebly het as daar nie langsaan 'n
pragtige kamp vol \VUgras{loodsoldaatjies} \textsuperscript{(29)} was wat my aandag
getrek het nie.

%%K t16-tit-2 sub
\VUsaut{4.0mm}
\VUpk{10.2pt}{40}{Soldate en oorlog.}
\VUsaut{2.8mm}

%%K t16-04-1 p
4. --- My oupa, wat \VUgras{invalide} \textsuperscript{(24)} is en wat die leër moes
verlaat weens 'n \VUgras{wond} waarvoor hy die \VUgras{militêre medalje} gekry het, hou
baie daarvan om vir my te vertel van die \VUgras{veldtogte} waaraan hy deelgeneem het.
Hy was gelukkig terwyl hy die verskillende bewegings van die klein leër in miniatuur
verduidelik het. 'n Groep egte soldate wat die basaar besoek het: 'n
\VUgras{geniesoldaat} \textsuperscript{(25)}, 'n \VUgras{Alpejagter}
\textsuperscript{(26)} met die kop bedek met 'n \VUgras{baret}, 'n
\VUgras{sersant-majoor} \textsuperscript{(27)} van die infanterie en 'n
\VUgras{artilleriesersant} \textsuperscript{(28)} met 'n goue \VUgras{tres} op die mou,
het naby ons stilgehou om na die verduidelikings te luister.

%%K t16-05-1 p
5. --- My oupa, heeltemal trots omdat hy so 'n blinkende gehoor gehad het, het 'n
volledige slagplan geïmproviseer.

%%K t16-05-2 p
Die versterkte stad, met sy \VUgras{walle} en \VUgras{gragte,} is beleër deur die
\VUgras{vyandelike leër,} wat ná 'n lang \VUgras{bombardement} gereed was om te
\VUgras{bestorm.} Maar die \VUgras{beleërdes,} deur die honger gedwing, het 'n
\VUgras{uitval} na die \VUgras{kampement} van die \VUgras{beleëraars} gewaag. Nadat
hulle die \VUgras{aanval} met 'n hardnekkige \VUgras{geveg} rondom die
\VUgras{tentekamp} afgeslaan het, het die \VUgras{oorwinnende} beleëraars hulle
\VUgras{eskadrons} uitgestuur om die \VUgras{oorwonne} aanvallers te agtervolg. Dié het
hulle \VUgras{teruggetrek} en die \VUgras{slagveld} met \VUgras{dooies} en
\VUgras{gewondes} bedek. \VUgras{Draagbaardraers} het laasgenoemdes na die
\VUgras{ambulans} weggedra om hulle deur die \VUgras{chirurg} te laat versorg.

%%K t16-05-3 p
Ondanks die heldhaftige verset van 'n paar \VUgras{bataljons} wat 'n \VUgras{vierkant}
gevorm het, was die \VUgras{nederlaag} volledig, die res van die leër het
\VUgras{gevlug} en baie \VUgras{gevangenes} agtergelaat. Die vyand het sy
\VUgras{oorwinning} gaan voltooi deur die stad te beset. Miskien sal dit die einde van
die veldtog wees en sal dit ná 'n \VUgras{wapenstilstand} moontlik wees om 'n
\VUgras{vredesverdrag} te teken.

%%K t16-tit-3 sub
\VUsaut{4.4mm}
\VUpk{10.2pt}{40}{Geskenke vir nuwejaar.}
\VUsaut{2.8mm}

%%K t16-06-1 p
6. --- Die jong soldate, wat gelag het terwyl hulle na die skilderagtige verhaal van die
veteraan geluister het, het hom nog 'n paar verduidelikings gevra. Wat my betref, ek het
om die winkel geloop om na 'n mooi \VUgras{kruisboog} \textsuperscript{(32)} en 'n
\VUgras{boog} \textsuperscript{(33)} met sy \VUgras{pyl} \textsuperscript{(31)} te kyk.
Daar het ek 'n ander versameling diere gevind: twee \VUgras{sangvoëls}
\textsuperscript{(34)}: 'n outomatiese \VUgras{nagtegaal} en 'n \VUgras{grasmossie} wat
uit volle keel gesing het, en 'n reeks vreemde diere: 'n \VUgras{vlermuis}
\textsuperscript{(35)}, 'n \VUgras{boa} \textsuperscript{(36)}, 'n \VUgras{skilpad}
\textsuperscript{(37)}, 'n \VUgras{rob} \textsuperscript{(38)}, 'n \VUgras{walvis}
\textsuperscript{(39)} en 'n \VUgras{haai} \textsuperscript{(40)} van vergulde leer.

%%K t16-07-1 p
7. --- Ek het nie lank stilgehou om na hulle te kyk nie, want ek het begin sien waarvan
ek gedroom het: 'n pragtige \VUgras{marionetteater} \textsuperscript{(41)} met
skitterende \VUgras{dekors} en 'n verruklike rooi \VUgras{gordyn.} Op die
\VUgras{toneel} het twee akteurs gelyk of hulle 'n \VUgras{rol} speel in 'n baie
boeiende \VUgras{toneelstuk,} want die klein kartonne \VUgras{toeskouers} wat in die
loges opgestel was of op die \VUgras{leunstoele} en in die \VUgras{parterre} agter die
\VUgras{orkesleier} gesit het, het lewendig geklap.

%%K t16-07-2 p
Jammer genoeg het daardie teater baie duur gekos.

%%K t16-08-1 p
8. --- Origens was dit moeilik om die speelgoed te kies, want in die uitstalvensters was
allerhande speelgoed opgestapel: 'n \VUgras{Harlekyn} \textsuperscript{(42)}, 'n
\VUgras{Pulcinella} \textsuperscript{(43)}, 'n \VUgras{Pierrot} \textsuperscript{(44)},
\VUgras{uile} \textsuperscript{(45)} wat met die vlerke slaan, fluitende \VUgras{merels}
\textsuperscript{(46)}, blaffende \VUgras{poedels}, 'n \VUgras{fonograaf}
\textsuperscript{(47)} en \VUgras{geduldspeletjies.} Een van daardie speelgoedstukke het
my baie vermaak: dit het 'n \VUgras{seeslag} \textsuperscript{(48)} voorgestel waarin 'n
\VUgras{skip} deur \VUgras{seerowers} aangeval is naby 'n land beplant met
\VUgras{klapperpalms.}

%%K t16-noto-1 noto
\VUnotes{143.2mm}{%
\textit{(*) Sien tabel nr. 6.}%
}

%%K t16-09-1 p
9. --- Ek kon al daardie wonders baie maklik bekyk, want net min mense was in die
warenhuis, skaars 'n paar slenteraars, 'n \VUgras{dragonder} \textsuperscript{(49)} en
'n \VUgras{invorderaar} \textsuperscript{(50)} wat 'n \VUgras{wissel} by die vernaamste \VUgras{kas}
\textsuperscript{(51)} aangebied het.

%%K t16-10-1 p
10. --- Ek het my oupa gevra waarvoor die boeke dien wat op planke agter die
\VUgras{kassierster} lê; hy het my geantwoord dat dit \VUgras{boekhouboeke} is.

%%K t16-tit-4 sub
\VUsaut{3.6mm}
\VUpk{10.2pt}{40}{Gapery om die winkel.}
\VUsaut{2.8mm}

%%K t16-11-1 p
11. --- Toe ons die speelgoedafdeling verlaat het, het ons na die saalmakery gegaan om 'n
blik te werp op die \VUgras{saals} \textsuperscript{(53)}, \VUgras{stiebeuels}
\textsuperscript{(54)}, \VUgras{tooms} \textsuperscript{(55)}, \VUgras{bitte}
\textsuperscript{(56)} en \VUgras{rysweepe.} Terwyl die \VUgras{afdelingshoof}
\textsuperscript{(57)} 'n paar inligtings aan my oupa gegee het, het ek die uitstalling
van die \VUgras{porselein en kristal} \textsuperscript{(58)} gaan bekyk waar mooi
\VUgras{slaaikomme} en fyn \VUgras{teepotte} \textsuperscript{(59)} is.

%%K t16-12-1 p
12. --- Om na die eerste verdieping te gaan, het ons agter die uitstalvenster van die
\VUgras{korsette} \textsuperscript{(60)} verbygegaan, naas die kousafdeling, waar baie
mooi \VUgras{onderbroeke} \textsuperscript{(63)} uitgestal was. Naby het 'n
\VUgras{winkeljuffrou} \textsuperscript{(61)} 'n \VUgras{koopster}
\textsuperscript{(62)} mooi sy \VUgras{weefsels} \textsuperscript{(64)} laat kies.

%%K t16-12-2 p
Terwyl ek met die trap opgeklim het, het ek opgemerk dat die warenhuis versier is met
Japanse \VUgras{lampions} \textsuperscript{(65)}, \VUgras{sonskerms}
\textsuperscript{(66)} en reusagtige \VUgras{palmbome} \textsuperscript{(70)}.

%%K t16-13-1 p
13. --- Op die eerste verdieping was daar nie veel om te sien nie; net meubels (*),
\VUgras{brandkaste} \textsuperscript{(67)}, \VUgras{laaikaste} \textsuperscript{(68)},
\VUgras{kinderwaentjies} \textsuperscript{(69)}. Maar die geheelbeeld van die warenhuis
was baie \VUgras{vermaaklik.}

%%K t16-14-1 p
14. --- Aan ons voete het ek die musiekinstrumente gesien: \VUgras{harpe}
\textsuperscript{(71)}, \VUgras{kitare} \textsuperscript{(72)}, \VUgras{mandoliene}
\textsuperscript{(73)}, \VUgras{klepkornette} \textsuperscript{(74)},
\VUgras{klarinette} \textsuperscript{(75)}, viole met hulle \VUgras{strykstok}
\textsuperscript{(76)} en selfs 'n \VUgras{grootrom} \textsuperscript{(77)}. Effens
verder, by die papierafdeling, het 'n \VUgras{student} \textsuperscript{(78)} by die
\VUgras{winkelbediende} \textsuperscript{(79)} oor 'n skrifsak afgeding. Naby hulle het
ek 'n mooi \VUgras{abc-boek} \textsuperscript{(80)} onderskei.

%%K t16-14-2 p
In die middel van die warenhuis is 'n hoë \VUgras{Kersboom} \textsuperscript{(81)}
opgerig wat heeltemal voorsien was van kersies, snuisterye en allerhande speelgoed:
\VUgras{houtperde} \textsuperscript{(82)}, \VUgras{poppe} \textsuperscript{(83)}, ens..

%%K t16-14-3 p
Die \VUgras{parfumerie} \textsuperscript{(84)} met sy \VUgras{tandmiddels} en
verskillende \VUgras{parfums} het 'n baie goeie geur versprei wat na ons toe uitgewasem
het.

%%K t16-tit-5 sub
\VUsaut{3.4mm}
\VUpk{10.2pt}{40}{Ons koop iets.}
\VUsaut{2.8mm}

%%K t16-15-1 p
15. --- Omdat my oupa die tyd te lank begin vind het, het ons met die groot trap
afgegaan. Hy het effens stilgehou voor 'n \VUgras{sterrekundige tabel}
\textsuperscript{(85)} wat, het hy vir my gesê, \VUgras{komete} \textsuperscript{(a)},
\VUgras{maan- en soneklipse} \textsuperscript{(b)}, 'n windroos met die \VUgras{vier
windstreke} \textsuperscript{(c)} \VUgras{(Noord, Suid, Oos en Wes)}, 'n kaart van die
twee \VUgras{halfronde} \textsuperscript{(d)}, die \VUgras{planete}
\textsuperscript{(e)} en die \VUgras{sonnestelsel} \textsuperscript{(f)} voorstel. Ek het
niks van dit alles verstaan nie, en ek is baie meer geboei deur die met tresse versierde
livrei van 'n \VUgras{loopjoggie} \textsuperscript{(86)} wat verbygegaan het, en deur die
mooi \VUgras{waaiers} \textsuperscript{(87)} wat naas my was, naby die \VUgras{beursies}
\textsuperscript{(88)} en die \VUgras{briewetasse} \textsuperscript{(89)}.

%%K t16-16-1 p
16. --- Ek het op die uitstalling ook \VUgras{vere} \textsuperscript{(90)} en
\VUgras{gordelgespe} \textsuperscript{(91)} bewonder, die mooi \VUgras{kunsblomme}
\textsuperscript{(94)} wat sierlik in mandjies gerangskik was. Die \VUgras{viooltjies},
\VUgras{violiere}, \VUgras{hiasinte} \textsuperscript{(95)}, \VUgras{geraniums}
\textsuperscript{(96)}, \VUgras{lelies} \textsuperscript{(97)} en \VUgras{kappertjies}
\textsuperscript{(98)} was baie goed nagemaak.

%%K t16-tit-6 sub
\VUsaut{4.4mm}
\VUpk{10.2pt}{40}{'n Geskenk van my oupa.}
\VUsaut{2.8mm}

%%K t16-17-1 p
17. --- Ek het my oupa gevra om vir my een van die \VUgras{tandeborsels}
\textsuperscript{(92)} te koop wat in 'n doos naas \VUgras{haarspelde}
\textsuperscript{(93)} was. Hy het toegestem en ek het self my aankoop by die kas gaan
betaal, waarnaas 'n belangrike \VUgras{besending} \textsuperscript{(100)} vir 'n
\VUgras{klant} \textsuperscript{(99)} voorberei is.

%%K t16-18-1 p
18. --- Skielik het 'n ligte opskudding ontstaan onder die mense wat naby die artistieke
\VUgras{bronswerk} \textsuperscript{(101)} stilgehou het. 'n \VUgras{Dief}
\textsuperscript{(102)} is pas op heterdaad deur 'n \VUgras{inspekteur}
\textsuperscript{(103)} betrap terwyl hy iemand probeer besteel het. 'n Mens het gesorg
dat 'n polisieman laat kom word om hom te \VUgras{arresteer,} en net toe ons uitgegaan
het, is die oortreder na die gevangenis gebring.
\VUsaut{6.0mm}
\VUfilet{22mm}
